\section{Analyse des exigences}

Le mandat demande la construction d'un agent généraliste de résolution de
problèmes. L'objectif n'est donc pas de programmer quatre solveurs indépendants,
mais de construire un moteur de recherche réutilisable pour plusieurs espaces
d'états.

Le \cref{tab:exigences} relie chaque exigence fonctionnelle à la décision de
conception adoptée dans \AUSSI{}.

\begin{table}[H]
\centering
\caption{Correspondance entre les exigences et la conception proposée.}
\label{tab:exigences}
\begin{tabularx}{\textwidth}{p{4cm}Y}
\toprule
\textbf{Exigence} & \textbf{Réponse de conception} \\
\midrule
BFS, UCS, Greedy et A* &
Les stratégies sont définies par une politique de priorité interchangeable. \\
\addlinespace

Tree-Search et Graph-Search &
Le moteur accepte un mode de recherche qui contrôle la gestion des états
visités et des chemins répétés. \\
\addlinespace

Moteur générique &
La fonction \code{SEARCH(problem, strategy, mode)} dépend seulement de
l'interface abstraite \code{Problem}. \\
\addlinespace

Quatre catégories de problèmes &
Chaque domaine implémente les opérations génériques : état initial, actions,
transition, test de but, coût et heuristique. \\
\addlinespace

Interface utilisateur &
Une interface en ligne de commande permet de choisir le problème, l'instance,
la stratégie, le mode et certaines limites de ressources. \\
\addlinespace

Statistiques &
Chaque exécution retourne un objet \code{SearchResult} contenant le statut,
la solution, le coût, la profondeur et les métriques d'exécution. \\
\addlinespace

Expérience comparative &
Un protocole de benchmark exécute plusieurs instances et exporte les résultats
dans un format tabulaire reproductible. \\
\bottomrule
\end{tabularx}
\end{table}

\subsection{Objectif de généralisation}

Le principe central de \AUSSI{} est de séparer les éléments invariants des
éléments spécifiques à un problème. Les éléments invariants sont le moteur de
recherche, la gestion de la frontière \code{OPEN}, la construction des noeuds,
la collecte des métriques et la reconstruction de solution.

Les éléments variables sont la représentation d'état, les actions disponibles,
la fonction de transition, le test de but, les coûts et l'heuristique. Cette
séparation permet d'ajouter un nouveau problème sans réécrire les algorithmes
de recherche.